\section{A uniform central Dirichlet second-jet sign under A}
\label{sec:analytic-central-dirichlet-sign}

Let $\chi$ be a nonprincipal real primitive Dirichlet character of
conductor $q$, and put $a\in\{0,1\}$ for its parity,
$\chi(-1)=(-1)^a$. Throughout this section the completion is
\begin{equation}\label{eq:analytic-dirichlet-completion}
 Q_\chi(s)=\left(\frac q\pi\right)^{(s+a)/2}
                \Gamma\left(\frac{s+a}{2}\right),\qquad
 \Lambda_\chi(s)=Q_\chi(s)L(s,\chi).
\end{equation}
This fixes every gamma and conductor factor used in the local
coefficient. The principal primitive character, whose completion
has poles, is excluded.

\begin{theorem}[NN-A032: positive central coefficient and exact side counts]
\label{thm:analytic-central-dirichlet-sign}
Let $m=\operatorname{ord}_{s=1/2}L(s,\chi)\ge0$. It is even.
The coefficient $C_\chi$ defined below satisfies, unconditionally,
$C_\chi>m+2-\pi/2$ for $a=0$ and
$C_\chi>m+\pi/2$ for $a=1$.
Under A, list the positive ordinates of all noncentral zeros
with their multiplicities. Then the exact identity is
\begin{equation}\label{eq:analytic-dirichlet-C-positive}
 \begin{split}
 C_\chi&:=2+\frac{L'}L(3/2,\chi)
                         +\frac{Q_\chi'}{Q_\chi}(1/2)\\
 &=m+\sum_{\gamma>0}\frac2{1+\gamma^2}
     +\begin{cases}2-\pi/2,&a=0,\\ \pi/2,&a=1.
       \end{cases}
 \end{split}
\end{equation}
The sum converges absolutely. In particular,
\[
 C_\chi\ge m+2-\pi/2>0\quad(a=0),\qquad
 C_\chi\ge m+\pi/2>0\quad(a=1).
\]

Assume A for the rest of the theorem. If $m>0$, write
$A_m=L^{(m)}(1/2,\chi)/m!$. Then $A_m>0$. For the shifted
Dirichlet deformation of NN-A004,
\[
 L_t(s,\chi)=\sum_{n\ge1}\frac{\chi(n)}{(n+t)^s},
\]
continued analytically in $s$ and locally in $t$, its $m$ zero
branches near the center have, for positive $t$, the expansions
\begin{equation}\label{eq:analytic-dirichlet-central-branches}
 s_j(t)=\frac12+c_jt^{1/m}
                  +\frac{C_\chi}{m}c_j^2t^{2/m}
                  +O(t^{3/m}),\qquad
 c_j^m=\frac{L(3/2,\chi)}{2A_m}>0.
\end{equation}
For all sufficiently small $t>0$, the exact counts are
\[
 \begin{array}{c|cc}
 &\Re s_j(t)<1/2&\Re s_j(t)>1/2\\ \hline
 m\equiv2\pmod4&m/2&m/2\\
 m\equiv0\pmod4,\ m\ge4&m/2+1&m/2-1.
 \end{array}
\]
In the second row the two branches with purely imaginary first
coefficient both bend strictly left at second order. No central
zero is asserted to exist; the side-count statement is
conditional on $m>0$. Neither quadratic central nonvanishing nor
simplicity follows from this theorem.
\end{theorem}

\begin{proof}
\emph{The exact functional-equation sign and entire completion.}
A nonprincipal real primitive character is quadratic. Its
quadratic Gauss sum has the sign
$\tau(\chi)=i^a\sqrt q$, so the root number in the convention
\eqref{eq:analytic-dirichlet-completion} is $+1$. A primary
locator, including the sign rather than just its square, is
B.~H.~Gross, \emph{Some remarks on signs in functional equations},
section~1, equations~(1.3)--(1.6), p.~2. There the sign is proved
by writing the completed quadratic $L$-function as the quotient
of the completed zeta functions of the quadratic field and
$\mathbb Q$. DLMF~25.15.5 gives the primitive functional equation
with the Gauss sum explicitly.

For completeness, the usual order-one analytic input can be
seen directly with these normalizations. Put
\[
 \theta_\chi(x)=\sum_{n\in\mathbb Z}n^a\chi(n)e^{-\pi n^2x/q},
 \qquad x>0,
\]
where the $n=0$ term is zero. Poisson summation in the residue
classes modulo $q$, using
\[
 \sum_{b\bmod q}\chi(b)e^{2\pi inb/q}=\tau(\chi)\chi(n),
\]
and the Gaussian Fourier transform gives
\[
 \theta_\chi(x)=\frac{(-i)^a\tau(\chi)}{\sqrt q}
        x^{-a-1/2}\theta_\chi(1/x)
              =x^{-a-1/2}\theta_\chi(1/x).
\]
For $a=1$ the needed transform follows by differentiating the
Gaussian transform once; this accounts for the factor $-i$.
Termwise Mellin integration for $\Re s>1$, followed by the
substitution $x\mapsto1/x$ on $(0,1)$, yields
\begin{equation}\label{eq:analytic-dirichlet-theta-entire}
 \Lambda_\chi(s)=\frac12\int_1^\infty\theta_\chi(x)
       \left(x^{(s+a)/2}+x^{(1-s+a)/2}\right)\frac{dx}{x}.
\end{equation}
The theta series decays exponentially for $x\ge1$; hence this
formula is entire, is invariant under $s\mapsto1-s$, and gives
$\log\max_{|s|\le R}|\Lambda_\chi(s)|=O_q(R\log(R+2))$ by
comparison with a gamma integral. Thus the entire function has
order at most one. Its zeros are precisely the nontrivial
$L$-zeros: the gamma poles are canceled by the trivial zeros,
with nonzero resulting values by the functional equation and
the Euler product on the opposite half-plane. At the endpoint
$s=0$ for an even character, the reflected point is $s=1$;
here one also uses Dirichlet's nonvanishing theorem
$L(1,\chi)\ne0$ (DLMF~25.15.9, with source Apostol, 1976,
p.~149), rather than the Euler product in its open half-plane.

\emph{Unconditional sign of the second-jet coefficient.}
Write $F(z)=\Lambda_\chi(1/2+z)$. All its zeros
$\alpha=\rho-1/2$ lie in $|\Re\alpha|\le1/2$ by the classical
critical-strip theorem, independently of A. The paired
Hadamard argument, exactly as in NN-A033, gives
\[
 \frac{F'(1)}{F(1)}
       =m+\sum_{\{\alpha,-\alpha\},\ \alpha\ne0}
                         \frac2{1-\alpha^2}.
\]
The series is absolutely convergent and is grouped under
conjugation. If $\alpha=b+ic$, then
$\Re(1-\alpha^2)=1-b^2+c^2>0$, whence
$\Re(1-\alpha^2)^{-1}>0$. Purely real and purely imaginary
pairs also contribute positive terms. There are noncentral
zeros, since the primitive $L$-function has infinitely many
zeros in the strip (DLMF~25.15(ii)). Hence $F'(1)/F(1)>m$.
The completion identity and the gamma differences calculated
below give, with no zero-location assumption,
\[
 C_\chi=\frac{F'(1)}{F(1)}+
       \begin{cases}2-\pi/2,&a=0,\\ \pi/2,&a=1.
       \end{cases}
\]
This proves the unconditional strict inequalities. It does not
fix the sign of the central leading coefficient without A.

\emph{Real positivity and the paired zero product under A.}
For real $\sigma>1$, every factor of the Euler product
\[
            L(\sigma,\chi)=\prod_p(1-\chi(p)p^{-\sigma})^{-1}
\]
is positive; this is DLMF~25.15.2. The gamma and conductor
factors are also positive. Assumption A excludes real zeros in
$(1/2,\infty)$, so continuity gives
$\Lambda_\chi(\sigma)>0$ throughout this interval. The
sign-$+1$ functional equation
makes $F$ even, so its central order $m$ is even. Its central
leading coefficient is positive, by its positivity for $z>0$.
Dividing by the positive number $Q_\chi(1/2)$ gives $A_m>0$
when $m>0$.

Under A all noncentral zeros of $F$ are $\pm i\gamma$. The
paired genus-one Hadamard product is
\begin{equation}\label{eq:analytic-dirichlet-positive-product}
 F(z)=Bz^m\prod_{\gamma>0}\left(1+\frac{z^2}{\gamma^2}\right),
                         \qquad B>0.
\end{equation}
Indeed the order-one factorization initially permits an
exponential $e^{b z}$; evenness forces $b=0$, and the paired
genus-one exponentials cancel. The sum of reciprocal squared
ordinates converges, so the product and its logarithmic
derivative converge normally on compact sets avoiding zeros.
At $z=1$ this gives the exact positive formula
\begin{equation}\label{eq:analytic-dirichlet-log-derivative}
       \frac{\Lambda_\chi'}{\Lambda_\chi}(3/2)
                    =m+\sum_{\gamma>0}\frac2{1+\gamma^2}.
\end{equation}

\emph{Cancellation of the conductor and the gamma differences.}
Since $L'/L=\Lambda_\chi'/\Lambda_\chi-Q_\chi'/Q_\chi$,
\[
 C_\chi=\frac{\Lambda_\chi'}{\Lambda_\chi}(3/2)
          +2+\frac12\left[
              \psi\left(\frac{1/2+a}{2}\right)
             -\psi\left(\frac{3/2+a}{2}\right)\right].
\]
The conductor terms have canceled exactly. The digamma
reflection and recurrence identities, DLMF~5.5.4 and~5.5.2,
give
\[
 \psi(1/4)-\psi(3/4)=-\pi,\qquad
 \psi(3/4)-\psi(5/4)=\pi-4.
\]
Substitution in \eqref{eq:analytic-dirichlet-log-derivative}
proves \eqref{eq:analytic-dirichlet-C-positive} and both strict
sign bounds. This uses the actual completion unit, without a
conductor-size estimate or a numerical verification.

\emph{The shifted Puiseux coefficient.}
Assume $m>0$, hence $m\ge2$. Write $x=s-1/2$ and
\[
 L(1/2+x,\chi)=A_mx^m+A_{m+1}x^{m+1}+O(x^{m+2}).
\]
Evenness of $\Lambda_\chi$ gives
$A_{m+1}/A_m=-Q_\chi'/Q_\chi(1/2)$. The shifted series of
NN-A004 is locally
\[
 L_t(s,\chi)=L(s,\chi)-sL(s+1,\chi)t+O(t^2).
\]
Put $t=u^m$ and $x=cu+du^2+O(u^3)$. The coefficient of $u^m$
gives $A_mc^m=L(3/2,\chi)/2$. Its $m$ roots are distinct.
The coefficient of $u^{m+1}$ then gives
\[
 mA_mc^{m-1}d+A_{m+1}c^{m+1}
       -\left(L(3/2,\chi)+\tfrac12L'(3/2,\chi)\right)c=0,
\]
or $d=(C_\chi/m)c^2$. Since $2m\ge m+2$, the $O(t^2)$ term
does not affect these two equations. The analytic implicit
function theorem after dividing by $u^m$ at each distinct
leading root gives all $m$ analytic branches in $u$, with the
stated error. This proves
\eqref{eq:analytic-dirichlet-central-branches}.

Finally put $R=(L(3/2,\chi)/(2A_m))^{1/m}>0$. The leading
coefficients are $c_j=Re^{2\pi ij/m}$. If $m\equiv2\pmod4$,
none is purely imaginary and exactly half have each sign of
their real part. If $m\equiv0\pmod4$, precisely two are
$\pm iR$; each has second coefficient
$-C_\chi R^2/m<0$. All other branches keep the sign of their
nonzero first real coefficient. Counting them gives the two
rows of the table, with no branch left on the line for
sufficiently small positive $t$.
\end{proof}

\paragraph{Geometric translation.}
The quadratic character is the rank-one Artin representation of
its quadratic finite etale cover away from the conductor; its
root number is the actual $+1$ in Gross's equations~(1.3)--(1.6).
At a hypothetical central zero of length $m$, the shifted
analytic divisor is finite flat of rank $m$ near the parameter
origin. The ramified base change $t=u^m$ separates its leading
sections. Complex conjugation pairs the two sections with
leading coefficients $\pm iR$, and the strictly negative real
second coefficient puts both on the left of the chosen real
normal line. This is a statement about the full local analytic
algebra and its completion unit, rather than a single label on
its reduced support. The finite-free discriminant/etale
criterion is SGA~1, Expos\'e~I, section~4, Proposition~10; no
arithmetic family realizing the shifted parameter, central
nonvanishing theorem, or positive integral cohomology is inferred.
